% q0051.tex — Nobody expects the Monetary Inquisition!
\question{
  id        = {0051},
  title     = {Nobody expects the Monetary Inquisition!},
  source    = {OBECON},
  year      = {2026},
  round     = {Seletiva IEO -- Phase C},
  language  = {English},
  topic     = {Macro},
  subtopic  = {Rational Expectations / New Classical Economics},
  type      = {objective},
  statement = {%
    A flexible-price economy expected a 20\% money supply increase; the actual
    increase was 40\%. According to New Classical theory with rational
    expectations, what happens to the price level and to output in the short
    run?
  },
  choices   = {%
    \item Both price level and output increase. Had agents expected zero
          increase, the output effect would be larger.
    \item Policy has no effect on output, only on the price level, since agents
          already expected expansion.
    \item Price level and output increase more than under zero expectations.
    \item Both rise; sticky prices cause the full 40\% increase to raise output
          gradually.
  },
  answer    = {A},
  solution  = {%
    Under New Classical rational expectations, only the \emph{unexpected}
    component of policy affects real output. Here the unexpected component is
    $40\% - 20\% = 20\%$, which raises aggregate demand and output in the
    short run alongside prices. Had agents expected 0\%, the entire 40\%
    expansion would be a surprise, producing a larger real output effect.
    Note that flexible prices rule out sticky-price mechanisms. (A) is correct.
  },
}
